\chapter{INTRODUÇÃO}
\label{cap:introducao}

O ensino de matemática enfrenta o desafio de atender estudantes com conhecimentos prévios, ritmos e dificuldades diferentes. Em turmas numerosas, o acompanhamento individual exige tempo, instrumentos adequados e interpretação pedagógica dos dados produzidos nas atividades. Técnicas computacionais podem apoiar esse processo ao organizar evidências de desempenho e tornar padrões mais visíveis, mas suas respostas não substituem a avaliação do professor nem representam, isoladamente, a aprendizagem do estudante.

Nesse contexto, recursos como aprendizado de máquina (\textbf{\textit{machine learning}}), analítica da aprendizagem (\textbf{\textit{learning analytics}}) e sistemas tutores inteligentes têm sido investigados para apoiar avaliações diagnósticas, estimar níveis de proficiência e orientar intervenções pedagógicas \cite{Machine2019_007, Implementation2025_000, Math2021_001}. O uso responsável dessas técnicas requer clareza sobre os dados empregados, os limites das inferências e a relação entre as saídas computacionais e os objetivos educacionais.

\section{Justificativa}

A relevância deste trabalho decorreu de duas necessidades complementares. A primeira consistiu em organizar a literatura, dispersa entre diferentes bases e áreas de pesquisa, sobre técnicas computacionais aplicadas ao ensino de matemática. A segunda consistiu em transformar as evidências encontradas em requisitos técnicos e pedagógicos para uma especificação de protótipo voltada ao apoio do professor.

Para professores, uma síntese crítica auxilia na compreensão das abordagens investigadas, dos resultados reportados e das limitações metodológicas recorrentes. Para estudantes e instituições, a pesquisa contribui ao explicitar que desempenho observado, proficiência estimada, competência e aprendizagem são conceitos relacionados, porém não equivalentes. Essa distinção reduz o risco de uma classificação algorítmica ser tratada como diagnóstico definitivo.

Do ponto de vista acadêmico, a revisão sistemática documentou as estratégias de busca, a seleção dos estudos e a avaliação de suas limitações metodológicas. O relato foi estruturado com apoio do PRISMA 2020 \cite{PRISMA2020}, enquanto os estudos incluídos foram examinados pelo MMAT 2018 \cite{MMAT2018}. A síntese da revisão e a fundamentação pedagógica sustentaram a definição dos requisitos, dos critérios de dados e modelagem, do protocolo de avaliação e da arquitetura de referência apresentados neste trabalho.

\section{Problema de Pesquisa}

Diante desse contexto, o trabalho respondeu ao seguinte problema de pesquisa:

Como identificar e sintetizar, por meio de revisão sistemática da literatura, as principais técnicas computacionais aplicadas ao ensino de matemática e converter essas evidências em uma especificação de protótipo que apoie o professor no diagnóstico de competências e dificuldades dos estudantes?

\section{Objetivos}

\subsection{Objetivo Geral}

Mapear e analisar sistematicamente as aplicações de técnicas computacionais no ensino de matemática e, a partir das evidências encontradas, elaborar uma especificação técnica e pedagógica de protótipo para apoio ao diagnóstico de competências.

\subsection{Objetivos Específicos}

\begin{itemize}
    \item \textbf{OE1}: realizar uma revisão sistemática da literatura, com relato apoiado pelo PRISMA 2020, sobre estudos publicados entre 2015 e 2025;
    \item \textbf{OE2}: identificar e categorizar as principais abordagens computacionais aplicadas ao ensino de matemática;
    \item \textbf{OE3}: classificar as finalidades pedagógicas atribuídas às aplicações encontradas;
    \item \textbf{OE4}: analisar criticamente as metodologias e limitações dos estudos incluídos;
    \item \textbf{OE5}: mapear lacunas técnicas, pedagógicas, metodológicas e éticas;
    \item \textbf{OE6}: manter um fluxo automatizado e auditável para coleta, processamento e exportação da revisão;
    \item \textbf{OE7}: derivar requisitos funcionais e não funcionais, critérios de seleção de dados e modelos, um protocolo de avaliação e uma arquitetura de referência para o protótipo.
\end{itemize}

\section{Delimitação do Trabalho}

O escopo executado compreendeu a revisão sistemática da literatura, a avaliação metodológica dos estudos incluídos e a elaboração de uma especificação conceitual de protótipo. Essa especificação consolidou princípios pedagógicos, requisitos, critérios para seleção de fontes de dados e modelos, procedimentos de avaliação e uma arquitetura de referência.

O desenvolvimento de uma aplicação funcional, o treinamento de modelos com uma base definitiva e a validação com participantes não integraram o escopo empírico desta pesquisa. Por essa razão, o trabalho não atribuiu ao protótipo métricas próprias de acurácia nem alegou eficácia pedagógica. Essa delimitação preservou a correspondência entre as afirmações apresentadas e os artefatos efetivamente produzidos.

\section{Estrutura do Trabalho}

O trabalho está organizado em sete capítulos. Após esta introdução, o Capítulo 2 apresenta a fundamentação teórica relacionada à aprendizagem matemática, à avaliação e às técnicas computacionais empregadas no contexto educacional. O Capítulo 3 descreve os procedimentos metodológicos adotados na revisão sistemática, incluindo a estratégia de busca, os critérios de seleção, a deduplicação dos registros, a avaliação metodológica e o procedimento empregado para derivar a especificação do protótipo. O Capítulo 4 apresenta os resultados da revisão sistemática, a síntese dos estudos incluídos, a avaliação pelo MMAT e as principais lacunas identificadas na literatura.

Com base nessas evidências, o Capítulo 5 apresenta a especificação conceitual do protótipo, contemplando princípios de projeto, requisitos, critérios para dados e modelagem, protocolo de avaliação, explicabilidade e arquitetura de referência. O Capítulo 6 integra os resultados da revisão, da avaliação metodológica, da fundamentação pedagógica e da especificação, discutindo suas implicações em relação ao problema de pesquisa. Por fim, o Capítulo 7 apresenta as conclusões, as limitações e as possibilidades de continuidade da pesquisa.

%--------- FIM INTRODUÇÃO------------
